\documentclass[11pt,letterpaper]{article}

\usepackage{graphicx}           % Graphics package
\usepackage{amsmath} 
\usepackage{mathrsfs}
\usepackage[colorlinks=true]{hyperref} 
\usepackage{color}
\usepackage{amsfonts}
\usepackage[textheight=9in, textwidth=7.5in, letterpaper]{geometry}
\usepackage[makeroom]{cancel}
\usepackage{cite}
\usepackage{sectsty}
\usepackage{empheq}
\usepackage{appendix}
\usepackage{enumerate} 
\usepackage{simpler-wick}
\usepackage{amssymb}

\sectionfont{\fontsize{12}{12}\selectfont}
\subsectionfont{\fontsize{12}{12}\selectfont}

\newcommand*\widefbox[1]{\fbox{\hspace{2em}#1\hspace{2em}}}
\newcommand{\LP}{\left(}
\newcommand{\RP}{\right)}
\newcommand{\mc}[1]{\mathcal{#1}}
\newcommand{\R}{\mathbb{R}}
\newcommand{\C}{\mathbb{C}}
\newcommand{\Z}{\mathbb{Z}}
\newcommand{\D}{\partial}
\newcommand{\KET}[1]{\left| #1 \right\rangle }
\newcommand{\BRA}[1]{\left\langle #1 \right| }
\newcommand{\IP}[2]{\left\langle #1 \middle| #2 \right\rangle}
\newcommand{\PA}[2]{\frac{\partial #1}{\partial #2}}



%%%%%%%%%%%%%%%%---------------------MOST USEFUL COMMAND EVER---------------------%%%%%%%%%%%%%%%%%%%%%%
\newcommand{\fixme}[1]{{\bf {\color{red}[#1]}}}
\newcommand{\draftmode}{\usepackage[notref,notcite]{showkeys}}
\providecommand*\showkeyslabelformat[1]{\normalfont\sffamily\footnotesize#1}

%%%%%%%%%%%%%%%


\setcounter{tocdepth}{4}
\setcounter{secnumdepth}{4}

\begin{document}

\title{Quantum Information}

\author{Mathew Calkins}

\date{\today}

 
\maketitle

\abstract{Extremely rough notes on quantum information. Maybe these will turn into something coherent, maybe not!}

\tableofcontents




\section{POVM Measurements}

\subsection{Pure States and Projective Measurements}
As typically introduced, the quantum mechanics of pure states $\KET{\psi}$ living in a Hilbert space $\mc{H}$ connects to experiments via projective measurements. Formally, a projective measurement is described by an indexed collection of projectors \begin{equation}
\{ \Pi_x: \mc{H} \to \mc{H} \ | \ x \in S \}, \quad S \subset \R
\end{equation}
which add up\footnote{Here ``addition" is meant generously to include integration.} to the identity \begin{equation}
\sum_{x \in S} \Pi_x = \mathbf{1}
\end{equation}
so that they equivalently specify an internal direct sum decomposition \begin{equation}
\mc{H} = \bigoplus_{x \in S} \mc{H}_s, \quad \mc{H}_s \equiv \Pi_x (\mc{H}).
\end{equation}
A projective measurement then randomly projects any prepared pure state $\KET{\psi} \in \mc{H}$ \begin{equation}
\KET{\psi} \xmapsto{\mbox{\tiny{proj}}} \frac{\Pi_x \KET{\psi}}{\lVert \Pi_x \KET{\psi} \rVert} = \frac{\Pi_x \KET{\psi}}{\sqrt{\BRA{\psi} \Pi_x \KET{\psi}}}
\end{equation}
and the corresponding $x \in S$ is called the measurement outcome. Assuming WLOG that $\KET{\psi}$ is properly normalized $\IP{\psi}{\psi} = 1$, the probability of any particular outcome is given by the Born rule \begin{equation}
\mbox{Prob}_x = \lVert \Pi_x \KET{\psi} \rVert ^2 = \BRA{\psi} \Pi_x \KET{\psi}.
\end{equation}
If $S \subset \R$ is anything more analytically interesting than a finite set, this object should be interpreted with care as a probability measure. For the present discussion, it suffices to imagine that $S$ is set.


\subsection{Entangled States and POVMs}
The above is all well and good, but the following simple and natural thought experiment suggests a more general kind of process which does not quite take the form of a projective measurement, but which intuitively warrants being called a ``measurement" of our system. Suppose we are interested in the physics of some Hilbert space $\mc{H}_A$, but that there exists a second space $\mc{H}_B$ and that we have the experimental power to perform general joint measurements. A general joint state \begin{equation}
\KET{\Psi}_{AB} \in \mc{H}_A \otimes \mc{H}_B
\end{equation}
is entangled, but for our purposes it will suffice to consider a pure state \begin{equation}
\KET{\Psi}_{AB} = \KET{\psi} \otimes \KET{\phi_{\mbox{\tiny{aux}}}}
\end{equation}
where we think of the second factor as an auxiliary state. Now suppose we perform a projective measurement on the joint state, described as above by a set of projectors \begin{equation}
\{ \Pi_{AB,x} \ | \ x \in S \subset \R\}.
\end{equation}
The Born rule applied to the joint state gives the probability of outcome $x \in S$ as \begin{equation}
\mbox{Prob}_x = \LP \BRA{\psi} \otimes \BRA{\phi_{\mbox{\tiny{aux}}}} \RP  \Pi_{AB,x} \LP \KET{\psi} \otimes \KET{\phi_{\mbox{\tiny{aux}}}} \RP.
\end{equation}
If we think of this as a function of $\KET{\psi}$ at fixed $\KET{\Phi_{\mbox{\tiny{aux}}}}$ and $\Pi_{AB,x}$, we can naturally rewrite the result as the diagonal matrix element of a quadratic form $Q_x$ on $\mc{H}$ with matrix elements \begin{equation}
Q_x: \mc{X}_A \times \mc{H}_A \to \C, \quad \BRA{\phi} Q_x \KET{\chi} \equiv \LP \BRA{\phi} \otimes \BRA{\phi_{\mbox{\tiny{aux}}}} \RP  \Pi_{AB,x} \LP \KET{\chi} \otimes \KET{\phi_{\mbox{\tiny{aux}}}} \RP.
\end{equation}
The RHS is antilinear in the first argument and linear in the second, so this defines a bona fide quadratic form, and we can identify \begin{equation}
\mbox{Prob}_x = \LP \BRA{\psi} \otimes \BRA{\phi_{\mbox{\tiny{aux}}}} \RP  \Pi_{AB,x} \LP \KET{\psi} \otimes \KET{\phi_{\mbox{\tiny{aux}}}} \RP = \BRA{\psi} Q_x \KET{\psi}.
\end{equation}
Recall that the polarization identity on quadratic forms \begin{equation}
Q(u, v) = \frac{Q(u+v, u+v) - Q(u,u) - Q(v,v)}{2}
\end{equation}
tells us a quadratic form's entire action is fixed by its diagonal values, so knowing the value of $\BRA{\psi} Q_x \KET{\psi}$ for all $\KET{\psi} \in \mc{H}_A$ (or more weakly by linearity, just its values on unit-normalized such states) fixes $Q_x$ uniquely.\\
\\
\
The defining properties of $\Pi_{AX,b}$ as a projective measurement, plus judicious use of the polarization identity, imply that each $Q_x$ is a Hermitian positive semi-definite operator\footnote{Here we are slightly redundant, as positive semi-definiteness is only defined for Hermitian operators to begin with.} \begin{equation}
Q_x = Q_x^\dagger, \quad Q_x \geq 0
\end{equation}
and that the collection taken together is a decomposition of the identity \begin{equation}
\sum_{x \in S} Q_x = \mathbf{1}_A.
\end{equation}
These two properties define a \textit{positive operator-valued measurement} (POVM). The reader can convince themselves that any such collection $\{Q_x\}_{x \in S}$ with these two properties can be cooked up as we did here, by a suitable choice of auxiliary state $\KET{\phi_{\mbox{\tiny{aux}}}} \in \mc{H}_B$ and joint projective measurement $\{\Pi_{AB,x} \}_{x \in S}$.


\section{Pure State Estimation, Symmetric Subspace}

In the words of Michael Walter, we new address the\\
\\
\
\textit{\textbf{Goal} (Pure State Estimation): Find a POVM $\{ Q_{\hat{\phi}} \}_{\hat{\phi} \in \Omega}$ on $(\C^d)^{\otimes n})$, with possible outcomes in the set \begin{equation}
\Omega = \left\{ \rho_{\hat{\phi}} = \KET{\hat{\phi}} \BRA{\hat{\phi}} \right\}
\end{equation} of pure states on $\C^d$, such that when we measure $\phi^{\otimes n} = \KET{\phi}^{\otimes n} \BRA{\phi}^{\otimes n}$, we obtain an outcome $\hat{\phi}$ that is ``close" to $\phi$. }\\
\\
\
As a small clarification: when we talk about outcomes, we are not saying that the POVM projects down to pure states on $\C^d$. Indeed, in general a POVM does not admit an interpretation as a projector with a pure state measurement outcome. We mean that we have some collection of pure states in $\C^d$; that we are using this collection as our indexing set for the POVM; and that when our measurement returns the outcome associated with $\KET{\hat{\phi}}$, we interpret this as our best guess for the state whose $n$-fold product is our initial pre-measurement state. A well-designed such POVM at variable $n$ should have the property that, in the limit of large $n$, this procedure should yield guesses $\KET{\hat{\phi}}$ which are extremely similar to $\KET{\phi}$ extremely often. This is plausible essentially because, given two unit-normalized states $\KET{\psi}, \KET{\psi} \in \C^d$, we have \begin{equation}
| \IP{\psi}{\xi}| < 1 \implies \left| \LP \BRA{\psi}^{\otimes n} \RP \LP \KET{\xi}^{\otimes n} \RP   \right| = | \IP{\psi}{\xi}|^n \xmapsto{n \to \infty} 0.
\end{equation}
So if two states are at all distinguishable, then their corresponding $n$-fold copies are very distinguishable for large enough $n$. Of course the order of predicates is important. This is a pointwise limit (for any two fixed distinct states, distinguishability goes to 1 as $n$ goes to $\infty$), not a uniform one (to the contrary, given any $n$ no matter how large, we could adversarially produce distinct states whose $n$fold products are arbitrarily hard to distinguish).


\subsection{Continuous POVMs and the Uniform Measure}
In order to get a typical guess $\KET{\hat{\phi}}$ which is arbitrarily close to the true state $\KET{\phi}$ at large $n$ no matter what the latter happens to be, our indexing set $\Omega \subset \C^d$ must at least contain states arbitrarily close to any possible true state. Rather than taking $\Omega$ itself to depend on $n$, we will take $\Omega$ from the state to be uncountably infinite, covering all distinct states in $\C^d$. As two states are distinct if they correspond to different rays, this means our indexing set is $\Omega \simeq \C \mathbb{P}^{d-1}$.\\
\\
\
Now the correct probabilistic interpretation of our outcomes must be more explicitly measure-theoretic. A continuous POVM is specified by an indexed collection of operators on $\mc{H}$ \begin{equation}
\{ Q_x : \mc{H} \to \mc{H} \}_{x \in \Omega}
\end{equation}
plus the auxiliary data of a measure $dx$ on $\Omega$.  We require as before that each operator be positive semi-definite and that the collection jointly be a decomposition of the identity \begin{equation}
Q_x \geq 0 \ \forall x \in \Omega, \mbox{ and } \int_\Omega dx Q_x = \mathbf{1}.
\end{equation}
For any $\psi \in \C^d$ and $x \in \Omega$, the probability density of outcomes \begin{equation}
p_\psi(x) = \BRA{\psi} Q_x \KET{\psi}
\end{equation}
gives rise to finite probabilities \begin{equation}
\mbox{Prob}_\psi (\mbox{outcome} \in S) = \int_S dx \BRA{\psi} Q_x \KET{\psi} \mbox{ for } S \subset \Omega \mbox{ measurable}
\end{equation}
and expected values \begin{equation}
\mbox{E}_\psi [f: \Omega \to \C] = \int_\Omega dx \BRA{\psi} Q_x \KET{\psi} f(x)
\end{equation}
So that all of the above makes analytic sense, the mapping $x \mapsto Q_x$ must be measurable.  Notice then that the mapping $Q$ which takes measurable subsets $X \subset \Omega$ to operators \begin{equation}
Q(X) := \int_X dx Q_x
\end{equation}
can itself be thought of as a measure on $\Omega$ taking positive(-semidefinite)-operator values, justifying the term POVM.\\
\\
\
For the particular problem of pure state estimation, even before we decide how to map states $\KET{\hat{\phi}}$ to operators $Q_{\hat{\phi}}$, we must decide on an appropriate measure. Conjugation by any fixed unitary as a symmetry of all observable physics, so any physically sensible measure must respect this symmetry. \begin{equation}
\int d \psi f(\psi) = \int d \psi f( U \psi U^\dagger).
\end{equation}
Just as there is a unique rotation-invariant measure on the sphere up to overall scaling, this symmetry uniquely fixed $d \psi$ to be the \textit{Haar measure} after we fix scaling by demanding that $d \psi$ be a probability measure.




\section{Classical Compression}
Suppose Alice has in hand a message of length $N$ consisting only of H's and T's, produced by repeatedly flipping a weighted coin which lands heads-up with probability $p$ and tails-up with probability $q = 1- p$. Alice wanted to send this message to Bob in some compressed way that doesn't actually require sending $N$ independent bits. This is possible if Alice and Bob both knows how this length-$N$ string was generate, in particular knows the value of $p$, as does Alice, and are willing to tolerate some small amount of error. We want to cook up a compression scheme. To evaluate such a scheme, we will only care about the per-message bits Alice must send Bob, ignoring any bits stored in writing down the compression and decompression dictionaries to begin with. This is sensible if, for example, Alice and Bob want to send a very large number of messages. We will also mostly only care about the large-$N$ performance. \\
\\
\
More concretely, we are going to pick a fix $\varepsilon$ independent of $N$, and restrict our compression protocol to messages which are \textit{$\varepsilon$-typical}, in the sense that the fraction of heads is within $\varepsilon$ of $p$. The number of such strings is smaller than $2^N$, so as long as Alice and Bob agree to a dictionary ahead of times, she can communicate such a message to him using a number of bits equal to the base 2 log of the number of such messages; for atypical messages, the compression scheme will simply send an error message. The probability of a random such string of length $N$ being $\varepsilon$-typical in this sense is \begin{equation}
\mbox{Prob} \LP \left| \frac{k}{n} - p \right| > \varepsilon \RP = \mc{O} \LP \frac{1}{N} \RP
\end{equation}
so at large $N$ the failure probability decays like $N^{-1}$. Not bad! This means that the fraction of random strings that our encoding scheme can handle goes to 1 as $N \to \infty$. The actual dictionary doesn't matter here; all that is left is to determine how many length-$N$ strings are $\varepsilon$-typical.




The trick here is that, while in principle there are $2^N$ distinct length-$N$ messages, at large $N$ nearly all random messages will be such that the fraction of heads is nearly $p$; we call such messages \textit{typical}. The number of strings of length $N$ with exactly $Np$ heads\footnote{Here we assume $Np$ is an integer; this isn't true for general $N,p$ of course, but for sufficiently large $N$ we can make it true by an arbitrary small rounding of $p$.} is \begin{equation}
\begin{pmatrix}
N  \\ Np
\end{pmatrix} = \frac{N!}{(Np)! (N(1-p))!}.
\end{equation}
From the Stirling approximation \begin{equation}
n! \sim \sqrt{2 \pi n} \LP \frac{n}{e} \RP^n
\end{equation}
we have roughly \begin{equation}
\begin{aligned}
\begin{pmatrix}
N  \\ Np
\end{pmatrix} & \sim \frac{\sqrt{2 \pi N} \LP \frac{N}{e} \RP^N}{\sqrt{2 \pi Np} \LP \frac{Np}{e} \RP^{Np} \sqrt{2 \pi N(1-p)} \LP \frac{N(1-p)}{e} \RP^{N(1-p)}}.
\end{aligned}
\end{equation}
Actually, it's easier to get what turns out to be an asymptotically tight upper bound. This is the sensible direction from which to bound, as it will in turn give an upper bound on the number of bits needed to compress a length-$N$ random string, which is equivalently a lower bound on the performance of our asymptotically lossless compression scheme. To get this easy upper bound, we take the binomial expansion \begin{equation}
\begin{aligned}
1 & = (x + (1-x))^n \\
& = \sum_{\ell = 0}^N \begin{pmatrix} N \\ \ell \end{pmatrix} x^\ell (1-x)^{N - \ell} \\
& \geq \begin{pmatrix} n \\ k \end{pmatrix} x^k (1-x)^{N - k}
\end{aligned}
\end{equation}
where in the last step we have chosen some $k \in \{0, 1, \ldots, n\}$ and assumed $x \in [0,1]$. Setting $x = k/N = p$ gives \begin{equation}
\begin{aligned}
\begin{pmatrix} N \\ k \end{pmatrix} & \leq x^{-k} (1-x)^{k-N} \\
& = \LP \frac{k}{N} \RP^{-k} \LP 1 - \frac{k}{N} \RP^{N - k} \\
& = 2^{-k \log_2 \LP \frac{k}{N} \RP - (N-k) \log_2 \LP 1 - \frac{k}{N} \RP } \\
& = 2^{N \left[ p \log_2 ( p ) - (N-k) \log_2 \LP 1 - p \RP \right]} \\
& \equiv 2^{-kN h_2(p)}
\end{aligned}
\end{equation}
where $h_2$ is the \textit{(binary) Shannon entropy function} \begin{equation}
h_2(p) \equiv - p \log_2 p - (1-p) \log_2 (1-p).
\end{equation}



\bibliography{bibfile}
\bibliographystyle{utphys}

\end{document}


